\section*{Streszczenie}

Hearthstone to jedna z najpopularniejszych cyfrowych kolekcjonerskich gier karcianych. Gracze podejmują w niej decyzje w warunkach niepełnej informacji, przy dużej liczbie możliwych akcji i wysokiej zmienności stanu gry. Te cechy sprawiają, że projektowanie skutecznych agentów komputerowych jest trudne i wymaga rozwiązania kilku problemów naraz --- od reprezentacji stanu planszy po wybór ruchu w ograniczonym czasie.

Konkursy, takie jak \textit{Hearthstone AI Competition}, popularyzują rozwój agentów dla Hearthstone oraz zapewniają wspólną platformę porównań i udostępniają gotowe implementacje prostych agentów. Mimo to wiele silnych rozwiązań konkursowych opiera się na złożonym przeszukiwaniu drzewa gry, co utrudnia ich bezpośrednie zastosowanie w prostszych, parametrycznych agentach optymalizowanych ewolucyjnie.

Jednym z podejść do tego problemu jest metoda ewolucyjna. García-Sánchez i in. zaproponowali agenta parametrycznego z ważoną funkcją oceny stanu planszy, trenowanego przez koewolucję konkurencyjną. Takie rozwiązanie ma jednak ograniczenia, ponieważ korzysta z prostej strategii ewolucyjnej, wąskiej 21-parametrowej funkcji oceny oraz klasycznego schematu koewolucji wewnątrz populacji.

W pracy zaproponowano pięć modyfikacji agenta i funkcji oceny, m.in. rozszerzenie przestrzeni wag do 28 i 63 parametrów, wariant ze znormalizowanymi cechami oraz fazowy podział oceny z płynnym przejściem między etapami gry. Dodatkowo podstawową strategię ewolucyjną zastąpiono adaptacyjną ewolucją różnicową SHADE w dwóch wariantach: koewolucyjnym i z mechanizmem Hall of Fame.

Eksperymenty przeprowadzono w środowisku SabberStone na ustalonym zbiorze talii. Wyniki pokazują, że zaproponowane modyfikacje osiągają wyższą skuteczność turniejową niż agent bazowy, a warianty SHADE przewyższają klasyczną koewolucję.

\textbf{Słowa kluczowe:} Hearthstone, sztuczna inteligencja w grach, algorytmy ewolucyjne, optymalizacja agentów.

\section*{Abstract}

Hearthstone is one of the most popular digital collectible card games. Players make decisions under incomplete information, with a large number of possible actions and high game-state variability. These properties make designing effective computer agents difficult and require solving several problems at once --- from board-state representation to move selection under a limited time budget.

Competitions such as the \textit{Hearthstone AI Competition} popularize agent development for Hearthstone by providing a common benchmarking platform and ready-made implementations of simple agents. However, many strong competition solutions rely on complex game-tree search, which makes them difficult to apply directly in simpler parametric agents optimized evolutionarily.

One approach to this problem is evolutionary learning. García-Sánchez et al. proposed a parametric agent with a weighted board-state evaluation function trained through competitive coevolution. However, this solution has limitations: it uses a simple evolutionary strategy, a narrow 21-parameter evaluation function, and a classical within-population coevolution scheme.

This thesis proposes five modifications to the agent and evaluation function, including extending the weight space to 28 and 63 parameters, a normalized-feature variant, and phased evaluation with smooth transitions between game stages. In addition, the basic evolutionary strategy was replaced with adaptive differential evolution SHADE in two variants: coevolutionary and with a Hall of Fame mechanism.

Experiments were conducted in the SabberStone environment on a fixed set of decks. The results show that the proposed modifications achieve higher tournament performance than the baseline agent, and SHADE variants outperform classical coevolution.

\textbf{Keywords:} Hearthstone, game artificial intelligence, evolutionary algorithms, agent optimization.
